\section{Related work}
\label{sec:related-work}

\subsection{Linking-based time-stamping}

The closest structural antecedent we found is the \emph{binary linking
scheme} of Buldas, Laud, Lipmaa, and Villemson \citep{buldas1998}, which
improves the linear hash-chain time-stamping of Haber and Stornetta
\citep{haber1991}. There, the $n$-th linking item is $L_n = H(H_n, L_{n-1},
L_{f(n)})$: each item hashes the previous item and one earlier item at a
power-of-two-structured distance $f(n)$, so that a verifying chain between
any two items has length $O(\log n)$ (at most $3\lceil \log_2 n\rceil - 2$),
and Buldas, Lipmaa, and Schoenmakers later gave optimally sized
time-stamps for such schemes using threaded authentication trees
\citep{buldas2000}. FlatMerkle's re-injection of $\CK[a]$ through
$\mathrm{MID}(a)$ plays the same role as the back-link $L_{f(n)}$: it is
what makes a distant point of the chain reachable in one step. The
differences are in what is hashed and what is proved. In a linking scheme,
each $L_n$ is the output of a separate hash invocation over several
inputs, and the object being authenticated is the one-way dependence
between two items, for relative temporal ordering. In FlatMerkle, the
back-links are compression-function chaining states folded into a single
Merkle--Damg\aa rd stream, so the commitment is the ordinary hash of one
byte string, and the statements proved are inclusion of an element in that
string and, via fraud proofs, malformation of its layout.

\subsection{Authenticated append-only skip lists}

Among authenticated data structures, the closest prior art we identified
is Maniatis and Baker's
\emph{authenticated append-only skip lists} (AASL)
\citep{maniatis2003}, itself building on the persistent authenticated
dictionary of Goodrich, Tamassia, and Schwerin
\citep{goodrich2001}. Both use the same deterministic,
power-of-two skip-index structure that FlatMerkle's position formula
(Lemma~\ref{lem:position}) also produces, and both achieve $O(\log N)$
membership and append proofs.

The mechanisms differ, however: AASL is a genuine skip list, with distinct,
separately tagged hash nodes at each level connected by explicit pointers,
computed via multiple separate hash-function invocations per node.
FlatMerkle achieves the equivalent skip structure without any separate
node objects or multiple hash invocations -- everything is a single,
continuous Merkle--Damg\aa rd stream, and the ``skip'' capability comes
entirely from re-injecting the same hash function's own internal state
back into that one stream. This is what allows FlatMerkle's digest to be a
literal, unmodified, single-pass hash of the raw structured bytes (Section
\ref{sec:bitcoin-annex}) -- a property we are not aware of AASL, its
predecessors, or the linking schemes above offering, since their
commitments are tied to their own multi-node hash structure rather than a
flat byte stream.

\subsection{Other related, but structurally distinct, work}

\begin{itemize}
\item \textbf{Proofs of sequential work}~\citep{mahmoody2013sequential} and
  related hash-chaining constructions also reuse intermediate hash states,
  but for a different purpose (proving elapsed sequential computation) and
  without the recursive, doubling checkpoint structure needed for
  logarithmic inclusion proofs.
\item \textbf{Historical/tamper-evident logging} (e.g.\ Crosby and Wallach's
  history trees \citep{crosby2009}) targets a related but distinct problem
  -- efficient proofs of consistency between two versions of an append-only
  log -- using an explicitly tree-shaped structure. The underlying
  probabilistic skip list data structure itself traces back to
  \citet{pugh1990}.
\item \textbf{Merkle mountain ranges}~\citep{todd2016mmr} generalize Merkle
  trees to efficiently-appendable forests of perfect binary trees, but
  remain tree-shaped structures, not a single linear hash stream.
\item \textbf{Bitcoin's own (unrelated) use of ``skip lists''} for block
  header lookups (e.g.\ for efficient median-time-past or difficulty
  retargeting queries) skips over \emph{block headers} via explicit stored
  pointers between blocks, not via hash-chain state re-injection, and
  solves a different problem (locating a specific historical block quickly)
  than proving inclusion inside a single committed dataset.
\end{itemize}

\subsection{Scope of this comparison}

This is a targeted literature search focused on the specific technique of
achieving logarithmic proofs from a single linear hash stream via
compression-state re-injection, not an exhaustive survey of authenticated
data structures or of the time-stamping literature, and not a patent
search. We are not aware of prior work using exactly this mechanism, but we note this as ``apparently
undocumented in what we found,'' not as a claim of novelty suitable for a
patent filing or a definitive related-work section without independent
verification.
